% ============================================================
%  吴宛幸  个人简历 (中文版 — 生物 / 健康方向)
% ============================================================
\documentclass[10pt, a4paper]{article}

% ── Encoding & CJK Support ──────────────────────────────────
\usepackage[UTF8, heading=false, fontset=none]{ctex}
\setCJKmainfont{Heiti SC}[AutoFakeBold=true, AutoFakeSlant=true]
\setCJKsansfont{Heiti SC}[AutoFakeBold=true]
\setCJKmonofont{Heiti SC}[AutoFakeBold=true]
\usepackage[T1]{fontenc}
\usepackage{palatino}
\usepackage{mathpazo}
\usepackage{microtype}

% ── Page Layout ──────────────────────────────────────────────
\usepackage[top=0.7cm, bottom=0.7cm, left=1.4cm, right=1.4cm]{geometry}
\setlength{\parindent}{0pt}
\setlength{\parskip}{0pt}
\linespread{0.96}\selectfont

% ── Packages ─────────────────────────────────────────────────
\usepackage{xcolor}
\usepackage{hyperref}
\usepackage{enumitem}
\usepackage{titlesec}
\usepackage{tabularx}
\usepackage{array}
\usepackage{fancyhdr}
\usepackage{tikz}
\usepackage{etoolbox}
\usepackage{graphicx}

% ── Colors ───────────────────────────────────────────────────
\definecolor{headcol}{RGB}{30, 30, 30}
\definecolor{rulecol}{RGB}{180, 50, 30}
\definecolor{secline}{RGB}{200, 200, 200}
\definecolor{linkcol}{RGB}{40, 80, 160}
\definecolor{dategray}{RGB}{100, 100, 100}

% Tag colors
\definecolor{tagroute}{RGB}{130, 50, 160}
\definecolor{taggenomics}{RGB}{40, 140, 90}
\definecolor{tagrl}{RGB}{180, 100, 20}
\definecolor{tagclin}{RGB}{30, 110, 160}
\definecolor{tagautoresearch}{RGB}{30, 100, 160}

\hypersetup{
  colorlinks = true,
  urlcolor   = linkcol,
  linkcolor  = linkcol,
  pdfauthor  = {吴宛幸},
  pdftitle   = {吴宛幸，个人简历},
}

% ── Section Headers ──────────────────────────────────────────
\titleformat{\section}
  {\large\bfseries\color{headcol}}
  {}{0em}{}
  [\vspace{2pt}{\color{secline}\rule{\linewidth}{0.4pt}}\vspace{-3pt}]
\titlespacing*{\section}{0pt}{2pt}{0pt}

% ── Badge Tags ───────────────────────────────────────────────
\newcommand{\badge}[2]{%
  \tikz[baseline=(x.base)]{
    \node[inner xsep=3.5pt, inner ysep=1.6pt, rounded corners=2.5pt,
          draw=#1, line width=0.45pt,
          text=#1, font=\fontsize{6}{7}\selectfont\sffamily]
      (x) {#2};
  }%
}

\newcommand{\troute}  {\badge{tagroute}{routing}}
\newcommand{\tgenomics}{\badge{taggenomics}{genomics}}
\newcommand{\trl}     {\badge{tagrl}{agentic}}
\newcommand{\tclin}   {\badge{tagclin}{clinical}}
\newcommand{\tautoresearch} {\badge{tagautoresearch}{autoresearch}}

% ── Bullet Lists ─────────────────────────────────────────────
\setlist[itemize,1]{
  leftmargin = 1.3em,
  label      = {\small\textbullet},
  itemsep    = 0pt,
  topsep     = 1pt,
  parsep     = 0pt,
}

% ── Helpers ──────────────────────────────────────────────────
\newcommand{\unilogo}[1]{%
  \raisebox{-0.5pt}{\includegraphics[height=8pt]{assets/logos/cv/#1}}%
}
\newcommand{\me}[1]{\textbf{#1}}
\newcommand{\eqstar}{$^{*}$}
\newcommand{\corr}{$^{\dagger}$}

\def\pubskip{1pt}

% Single publication block
\newcommand{\pub}[4]{%
  \vspace{\pubskip}%
  \noindent\textbf{#1}\ifstrempty{#2}{}{\enskip #2}%
  \begin{itemize}[topsep=0pt]
    \item #3
    \item \textit{#4}
  \end{itemize}%
}

% Publication block with extra "highlight" bullet
\newcommand{\pubh}[5]{%
  \vspace{\pubskip}%
  \noindent\textbf{#1}\ifstrempty{#2}{}{\enskip #2}%
  \begin{itemize}[topsep=0pt]
    \item #3
    \item \textit{#4}
    \item #5
  \end{itemize}%
}

% Education / Experience entry
\newcommand{\entry}[2]{%
  \noindent
  \begin{tabularx}{\linewidth}{@{}X r@{}}
    \textbf{#1} & {\small\color{dategray}#2}
  \end{tabularx}%
  \par\vspace{0pt}%
}

% ── Header / Footer ──────────────────────────────────────────
\pagestyle{fancy}
\fancyhf{}
\renewcommand{\headrulewidth}{0pt}
\fancyfoot[R]{\footnotesize\color{dategray}\thepage}

% ============================================================
\begin{document}
\thispagestyle{empty}

% ── NAME BLOCK ───────────────────────────────────────────────
\noindent
\begin{tabularx}{\linewidth}{@{}X r@{}}
  \raisebox{-4pt}{%
    {\fontsize{28}{32}\selectfont\bfseries 吴宛幸}%
  }
  &
  \begin{tabular}[b]{@{}r@{}}
    \small 150 6798 2279 \quad \href{mailto:wuwanxing3574@gmail.com}{wuwanxing3574@gmail.com}\\[2pt]
    \small\href{https://github.com/wwx777}{GitHub}
  \end{tabular}
\end{tabularx}

\vspace{4pt}
{\color{rulecol}\rule{\linewidth}{1.2pt}}
\vspace{1pt}

% ── ABOUT ────────────────────────────────────────────────────
\section{个人简介}

\textbf{巴黎理工学院（IP Paris）}计算机科学硕士（Digital Skills for Health Transformation track）在读，本科毕业于\textbf{南方科技大学}（数据科学与大数据技术），现为 \textbf{SUSTech-NLP} 科研助理（导师：陈冠华教授）。我的研究方向是\textbf{异构能力的 LLM 应该如何协作}。在 \textbf{RouterXBench}（per-query 路由）中，我设计了一套将 router 内在能力与场景适配性解耦的三轴评估框架，并用它发现\emph{模型内部 hidden state} 比 verbose、logit、外部 embedding 等其他信号都更适合做 routing。在 \textbf{Scope-then-Cover}（agent 适应）中，我利用 exploration 内部的能力非对称：强 actor 提名一个高 recall 的小子空间，冻结的廉价 worker 并行覆盖，封装为单一 \texttt{explore} action，可挂入 test-time adaptation (TTA) 与 online experiential learning (OEL) 两种 host。

% ── EDUCATION ────────────────────────────────────────────────
\section{教育经历}

\entry{\unilogo{ipp.png}\enskip 巴黎理工学院（Institut Polytechnique de Paris）}{法国巴黎}
\textit{计算机科学硕士 \textnormal{\small（Digital Skills for Health Transformation track）}} \hfill \textit{2025.09 -- 至今}
\begin{itemize}
  \item \textbf{ReflexSQL}：基于探索与记忆的 Text-to-SQL agent，验证于 Spider~2.0-Lite \quad [\href{https://wwx777.github.io/assets/reflexsql-poster.pdf}{海报}]。
\end{itemize}

\vspace{1pt}
\entry{\unilogo{sustech.png}\enskip 南方科技大学}{深圳}
\textit{数据科学与大数据技术，理学学士} \hfill \textit{2021.09 -- 2025.07}
\begin{itemize}
  \item 本科毕业论文：大小模型协同路由（在 \textbf{RouterXBench} 中延续）。
\end{itemize}

% ── PUBLICATIONS ─────────────────────────────────────────────
\section{学术论文}

\pubh{Towards Fair and Comprehensive Evaluation of Routers in Collaborative LLM Systems}{\troute}
    {\me{Wanxing Wu}, He Zhu, Yixia Li, \ldots, Guanhua Chen\corr}
    {在审，EMNLP 2026 \quad[\href{https://arxiv.org/abs/2602.11877}{arXiv:2602.11877}]}
    {提出 \textbf{RouterXBench}：3 轴评估框架（AUROC / LPM-MPM-HCR / 6 个 ID-OOD benchmark）；用它系统验证了 \textbf{模型内部 hidden state 是比 verbose / logit / 外部 embedding 更具判别力的 routing 信号}，结论在不同族、不同规模的小模型上均成立。}

\pubh{Can Weak Models Explore? Scope Then Cover for LLM Agent Adaptation}{\trl}
    {\me{Wanxing Wu}\eqstar, He Zhu\eqstar, Guanhua Chen\corr \quad ($^*$同等贡献)}
    {在审，2026}
    {LLM agent 适应新环境的关键在于 exploration。我们用两个受控 probe 观察到两类现象：当环境对目标位置无信息时，scale up 探索者得不到任何增益；当候选空间超出单步 inspection budget 时，各种规模都退化到 random search。两者揭示了 exploration 内部的同一种结构：定位一个有希望的子空间是 capability-bound 的，而验证目标是否落在其中并不是。为此提出 \textbf{scope-then-cover}：capable actor 提名一个小但高 recall 的子空间，冻结的廉价 worker 并行覆盖；整体封装为单一 \texttt{explore(scope, query)} action。同一个 action 既能挂在冻结的 test-time actor 上，也能挂在 online experiential learning host 上，在 ALFWorld、WebShop、SWE-Bench-Verified 上都超过 no-explore baseline。}

\pubh{AutoInterpBench: Level-1 Grounded Benchmarking for Autonomous Mechanistic Interpretability Discovery}{\tautoresearch}
    {Xuanzhe Xu, He Zhu, \me{Wanxing Wu}, Guanhua Chen\corr}
    {在审，EMNLP 2026}
    {提出 \textbf{AutoInterpBench}：面向\emph{自动化研究 (autoresearch)} agent 的三层基准（264 任务，764 个 GPU 实验，14 个模型），配套 \textbf{harness-driven discovery} 闭环，以 L1 分数作 fitness 信号迭代优化 agent 的 MI analysis primitives。}

% ── RESEARCH EXPERIENCE ──────────────────────────────────────
\section{科研经历}

\entry{\unilogo{sustech.png}\enskip SUSTech-NLP 实验室}{深圳}
\textit{科研助理}，导师：陈冠华教授 \hfill \textit{2025.01 -- 至今}

上述两篇论文的（共同）一作。RouterXBench/XRouter 中，基于 vLLM 搭建多模型推理与评测 pipeline，覆盖 batch inference、router evaluation，以及 logit/embedding/hidden-state 路由基线；Explore-Agent 中，使用 OPCD、skill-sd 与 Agent-RL 训练 Explore-Agent，部分成果已形成在投论文 Scope-then-Cover。

\vspace{1pt}
\entry{\raisebox{-0.5pt}{\includegraphics[height=8pt, width=16pt, keepaspectratio]{assets/logos/cv/bgi.png}}\enskip 华大基因（BGI）智能医学研究院（IIMR）}{深圳}
\textit{机器学习实习生（AI for Genomics）} \hfill \textit{2024.07 -- 2025.01}

搭建 cfDNA 数据处理与机器学习建模 pipeline，用于结直肠癌早筛 \tgenomics，结合 cfDNA 表格特征与 BERT-style end-motif 序列编码器；负责特征清洗、统计筛选、模型导向特征选择、集成建模与交叉验证。在 held-out 队列上达到 \textbf{94.88\%} 灵敏度。

\vspace{1pt}
\entry{\unilogo{pku.png}\enskip 北京大学深圳医院}{深圳}
\textit{数据分析实习生} \hfill \textit{2023.07 -- 2023.10}

\textbf{临床结果分析} \tclin：对多时间点临床记录做生物统计分析，比较不同术式之间的术后结果。

% ── HONORS ───────────────────────────────────────────────────
\section{荣誉与奖项}

\begin{itemize}[itemsep=0pt, topsep=1pt]
  \item \textbf{优秀毕业生}，南方科技大学致仁书院（2025.07）
  \item \textbf{优秀毕业论文}，Top 5\%，南方科技大学（2025.07）
  \item 优秀学生奖学金，南方科技大学（2024.10）
\end{itemize}

% ── SKILLS ───────────────────────────────────────────────────
\section{技能}

\noindent\textbf{大模型与机器学习：} 推理框架（vLLM、SGLang），Agentic-RL 工作流。 \quad \textbf{编程与工具：} Python、PyTorch、Docker、Linux、Git。\\
\noindent\textbf{语言：} 英语（IELTS 6.5，CET-4，CET-6）；中文（母语）。

\end{document}
